\thispagestyle{empty} %%<- poner en mas paginas si los resumenes ocupan varias (elimina la numeracion de la pagina actual)
\pagebreak


\begin{flushleft} \hyphenpenalty=10000\exhyphenpenalty=10000 {\bf $\bullet$ Resumo:\;\;}
%
Aqu\'{\i} vai o resumo en galego. Debe coincidir co introducido na secretar\'{\i}a virtual no momento do dep\'osito da memoria final e solicitude de defensa. Dado que a aplicaci\'on inform\'atica de secretar\'{\i}a virtual non admite calqueira caracter, o regulamento permite introducir nela representaci\'ons alternativas dos caracteres problem\'aticos (por exemplo introducir gamma en vez de $\gamma$, introducir YBa2Cu3O(7-delta) en vez de \mbox{YBa$_2$Cu$_3$O$_{7-\delta}$}, etc.). Non ten por qu\'e coincidir co resumo do proxecto de TFG/TFM que foi feito no momento da asignaci\'on de traballo e titor (entendendo que o titor da o visto bo no seu informe final). Ten que haber resumos en (como m\'{\i}nimo) galego, castel\'an e ingl\'es, cada un de 300 palabras m\'aximo.
%
\end{flushleft}\mbox{}


\begin{flushleft} \hyphenpenalty=10000\exhyphenpenalty=10000 {\bf $\bullet$ Resumen:\;\;}
%
Aqu\'{\i} va el resumen en castellano. El orden de los idiomas puede cambiarse a voluntad. Tambi\'en (y aunque la hoja de res\'umenes no contabiliza para el n\'umero l\'{\i}mite de p\'aginas, ver secci\'on \ref{quecontabiliza} de este documento) puede  reducirse, si se desea, el tama\~no de letra de los res\'umenes en los dos idiomas que no sean el usado  en el texto principal de la memoria (por ejemplo, anteponiendo \verb \footnotesize{}  al texto). % anteponiendo \footnotesize{} al texto, por ejemplo
%
\end{flushleft}\mbox{}


\begin{flushleft} \hyphenpenalty=10000\exhyphenpenalty=10000 {\bf $\bullet$ Abstract:\;\;}
%
The abstract goes here.
%
\end{flushleft}\mbox{}

\begin{center}
\rule{61mm}{0.1mm}\\
\end{center}
{\bf Agradecementos (opcional)}
Aqu� p�dense incluir os agradecementos, se non foron inclu\'{\i}dos xa na portada interior.
